\section{Discussion}

The central finding of this audit is a stark dissociation between in-distribution competence and cross-dataset generalization. Within SIPaKMeD, the distilled EfficientNet-B0 student is an accurate, well-calibrated, and visually faithful five-class classifier: it reaches an accuracy of $0.966$, a macro-F1 of $0.967$, a one-vs-rest AUROC of $0.998$, and an Expected Calibration Error of only $0.015$, with Grad-CAM attributions consistently localized to nuclear and cytoplasmic structure rather than background artefacts. Yet when this same model is confronted with three independent Papanicolaou datasets under a deterministic binary Normal/Abnormal alignment, performance collapses toward, and in places below, majority-class baselines: binary accuracy falls to $0.639$ on Herlev, $0.605$ on HiCervix, and $0.532$ on APACC, the latter essentially indistinguishable from chance (AUROC $0.525$). The roughly thirty-three-point internal-to-external gap is the empirical heart of this paper, and it directly answers RQ2 in the negative: strong within-dataset accuracy does not translate into deployable cross-dataset performance.

It is important to read the student-exceeds-teacher result in the correct light. The B0 student surpasses its B4 teacher on every in-distribution metric---accuracy ($0.966$ versus $0.943$), balanced accuracy, macro-F1, kappa, AUROC, and calibration---at roughly a quarter of the parameters ($4.67$\,M versus $18.5$\,M). Consistent with the regularizing view of knowledge distillation \citep{hinton2015distilling}, the soft-target signal appears to smooth the student's decision boundaries and improve generalization \emph{within the training distribution}. This is a genuine efficiency win and confirms H1. However, the same audit shows that this advantage is confined to the source domain: a better in-distribution model is not a more transportable one. Distillation optimizes agreement with a teacher trained on SIPaKMeD, and any domain-specific shortcuts \citep{geirhos2020shortcut} the teacher encodes are inherited, not corrected, by the student.

The diagnosis isolates the failure mechanism and, crucially, exonerates the label crosswalk. Because the class-conditional separation of $P(\text{abnormal})$ is positive in all three datasets ($+0.284$ Herlev, $+0.112$ HiCervix, $+0.045$ APACC) and every AUROC exceeds $0.5$, the mappings cannot be inverted or mis-specified---a flipped alignment would drive AUROC below $0.5$. The model therefore retains \emph{some} valid ranking signal everywhere; what it loses is the operating point. The exemplary in-distribution calibration does not transfer: the baseline of $P(\text{abnormal})$ drifts per domain, sitting high on APACC (the model over-calls abnormal) and low on HiCervix (it under-calls), so a fixed threshold $\tau=0.5$ throws away the surviving ranking and pushes accuracy beneath the majority baseline. The oracle best-threshold analysis quantifies this calibration gap precisely---Herlev recovers to $0.75$ at $\tau=0.05$, HiCervix to $0.65$ at $\tau=0.85$---and confirms that miscalibration, not loss of discriminative information, is the dominant culprit. This resolves RQ3 in favour of H3: the collapse is genuine domain shift compounded by operating-point miscalibration, not a relabelling artefact. The ordering of degradation by slide-preparation distance---conventional Herlev $>$ liquid-based HiCervix $>$ auto-segmented APACC---further implicates stain and preparation domain shift as the underlying axis.

For screening deployment the implications are sobering and concrete. A model validated only within its source dataset, however impressive its headline numbers, can silently degrade to chance on the next site's slides, and a fixed decision threshold is actively harmful because calibration is the first thing to break. The remedies follow directly from the diagnosis. First, evaluation must be leakage-safe: same-field-of-view cells should be grouped so that no cluster straddles the train/validation boundary, since the random per-cell split renders even the $0.966$ internal figure optimistic and widens the true gap. Second, stain and domain normalization at input \citep{macenko2009stain} should be used to shrink the preparation distance that drives the degradation ordering. Third, because the ranking signal survives but the operating point does not, per-domain recalibration of the threshold or of the probabilities themselves \citep{guo2017calibration} is essential before any fixed cut-off is trusted on a new site. Finally, reporting practice should favour threshold-free AUROC alongside the excluded fraction, so that discriminative ability and label-space honesty are both visible rather than masked by a single, non-transportable accuracy.

\subsection{Limitations}

Several limitations bound these conclusions. First, the model is trained on a single source dataset (SIPaKMeD); the failure modes we characterize are therefore relative to one conventional-smear domain, and we cannot disentangle how much of the collapse a multi-source training corpus would absorb. Second, the in-distribution metrics rest on a random per-cell $80/20$ split that leaks same-field-of-view cells across train and validation; the reported $0.966$ is consequently optimistic, the genuine internal-to-external gap is larger than the headline thirty-three points, and we quantify the direction but not the exact magnitude of this leakage. Third, the cross-dataset evaluation collapses rich dysplasia and Bethesda gradings into a binary Normal/Abnormal axis---the only label level that is deterministically and many-to-one mappable across schemes---which necessarily discards high-grade detail: severity gradients, glandular lesions (AGC, Herlev columnar), and microbial categories are either excluded or flattened, so the binary numbers cannot speak to high-grade lesion detection specifically. Fourth, Grad-CAM evidence is attribution rather than proof of clinical reasoning, and the qualitative panel covers one representative cell per class. Finally, this is an algorithmic audit on retrospective public datasets with no prospective clinical validation, no reader study, and no assessment on prospectively collected slides; the remedies we propose are evidence-led hypotheses for closing the gap, not validated deployment recipes.

\section{Conclusion}

We have presented a leakage-aware, cross-dataset audit of a knowledge-distilled cervical-cytology classifier and used it to separate genuine model failure from labelling artefact. Answering RQ1, a $4.67$\,M-parameter EfficientNet-B0 student, distilled from an $18.5$\,M EfficientNet-B4 teacher, matches and in fact exceeds the teacher on every in-distribution metric (accuracy $0.966$ versus $0.943$) while remaining well-calibrated (ECE $0.015$) and explainable, confirming that distillation can deliver an efficient, faithful source-domain model. Answering RQ2 and RQ4, this competence does not transport: binary accuracy falls to $0.639$, $0.605$, and $0.532$ across Herlev, HiCervix, and APACC, the in-distribution calibration fails to transfer, and the degradation tracks slide-preparation distance from the conventional-smear source. Answering RQ3, the collapse is not a label-mapping artefact---positive class-conditional separation and AUROC above $0.5$ in every dataset, together with oracle-threshold recovery, identify genuine domain shift and operating-point miscalibration as the cause. The practical message is that within-dataset accuracy is a poor proxy for screening readiness, and that leakage-safe evaluation, stain and domain normalization \citep{macenko2009stain}, per-domain calibration \citep{guo2017calibration}, and threshold-free reporting with the excluded fraction are the necessary next steps toward cervical-cytology models that generalize beyond the dataset on which they were born.